\documentclass[12pt,letterpaper]{article}
\setlength{\textwidth}{6.5in}
\setlength{\oddsidemargin}{0in}
\setlength{\textheight}{9in}
\setlength{\topmargin}{-.5in}
\setlength{\parindent}{1.5em}
\setlength{\parskip}{.65\baselineskip}
\setcounter{secnumdepth}{0}
\setcounter{tocdepth}{1}
\newcommand{\entry}[1]{\clearpage\section{#1}}
\begin{document}
\begin{titlepage}\centering\vspace*{2in}
{\Large\bfseries Class Journal, Part II\par}
Avi Chetlin\par Department of Mathematics and Department of Physics\par Case Western Reserve University\par
ORBH 375: Building Leadership Character\par Dr. Susan Case\par
18 November 2026\end{titlepage}
\tableofcontents
\entry{Week 8. Authentic Leadership \textit{(October 13 and 15, 2026)}}
\subsection{Concrete Experience}
At the beginning of the semester I wrote that I had very little desire to be the Guy. As it turns out, I'll be leading the math club next year. My first impulse in trying to reconcile these two facts has been to worry that I'm taking up too much space. I already teach calculus, do research in the department, and spend a substantial amount of my time doing mathematics. The prospect of having yet another occasion on which people will hear from me gives me pause. At the same time, I do want to do the job. I like these people, I care about mathematics, and I would like the club to be somewhere people can enjoy doing it together.
The club matters a lot to me, and I feel as though I have to step up, or nobody else will.
Jonah Reinis (the current president) and I have done a lot of work to build our membership, and I don't want to see that go to waste.

\subsection{Reflective Observation}
What, exactly, do I mean by taking up too much space? There is a perfectly intelligible objection to a club in which one member gets to decide what is interesting and everybody else gets to attend. I would object to that club even if I happened to share that member's interests. But my worry has an odd feature. I feel guilty now, before I've led a single meeting in this capacity, and before I could possibly know whose idea has gone unheard. Apparently the mere fact of being visible is enough for me to feel as though I've already deprived somebody of something. I think I owe the other members a little more credit than that. They might have perfectly good uses for somebody willing to take responsibility for organizing the things they want to do.

\subsection{Abstract Conceptualization and Reading Reflection}
The Society of Leadership Fellows (2019a, pp. 8--9) asks members of a leadership team to describe the self they would most like their colleagues to encounter. I like the question because it exposes a difficulty in the advice to be authentic. My reluctance to take up space is quite sincere. So is my wish to run a club that someone new to mathematics might feel comfortable joining. Which of these impulses has the better claim to represent me? I would like a member to be able to say of me that when they suggested something, I helped them work out how to do it. That answer tells me considerably more about how to conduct a meeting than an inventory of my sincere feelings does.
Its moreover extremely important to me that peoples' takeaway from club meetings is not, ``wow, that guy sure likes to talk,'' but that they felt welcomed and comfortable in the environment.
I sometimes find it difficult to straddle the line between taking control to create the environment I want, and step back enough to let people feel unconstrained.

The October 15 prompts ask what respect meant to King and Gorman, and what I need in order to feel respected. King (1963) asks his audience to imagine how his four children will be judged when they grow up. Gorman (2021) invites us to look at the person already standing in front of us, a young Black woman who could once imagine becoming president and has now been asked to read a poem for one. I have always found King's repetition moving, because each return to the dream asks his listeners to picture yet another consequence of taking equality seriously. Gorman makes a related demand of the audience, simply by telling us who is speaking.
In my own experience, I feel respected when someone takes enough interest in an idea of mine to ask me how I arrived at it. I can do that for the person who comes to math club with a question that seems elementary to me. Asking what led them to it gives them a chance to teach me something about why the question is interesting in the first place.
And, I can remember how I felt when I was an underclassman with less of a background in math.

\subsection{Assessment Reflection}
My responses on the twelve-item authenticity scale, in order, were 5, 3, 4, 4, 5, 5, 3, 3, 4, 3, 4, and 2. Using the definitions in Newman's chapter on authenticity (2019, ch. 4), I scored 16 for authentic living, 18 for accepting external influence, and 11 for self-alienation. Authentic living concerns the extent to which I act according to my values and beliefs across the situations I ordinarily encounter. Accepting external influence includes both how readily others influence me and how strongly I feel obliged to meet their expectations. Self-alienation concerns the gap between what I am experiencing and what I consciously recognize about that experience. I find the last distinction particularly interesting. I can give somebody an entirely sincere account of how I'm doing and learn, some weeks later, that there were things about my own experience I hadn't understood well enough to tell them. That possibility gives me a reason to take the 11 seriously without imagining that a questionnaire could tell me exactly what I've overlooked.
I also think its interesting that the questionnaire seems to consider accepting externalal influence a bad thing.
We discussed this in class, but I think that at my current stage in life, it is often a good thing for me to accept external influence, given my lack of experience and the number of people with more experience surrounding me.

When I think about how I might answer as president of the math club, one difficulty is deciding what it means for somebody else's opinion to influence me. If a first-year student tells me that all our talks assume mathematics they haven't taken, I sincerely hope their opinion will influence me. If I choose a topic mainly because I think a professor will be impressed by it, I'd have a harder time defending the same willingness to be influenced. In the former case I can explain why I've changed my mind by reference to the kind of club I want to run. In the latter, I suspect I'd have some trouble explaining my decision to the student who came hoping to learn something.

The True North Mirror gives me a second point of comparison. For self-awareness I scored 62 against an average of 72 among leaders from around the world. For values the comparison was 71 to 79, for integrated life 39 to 61, for support team 58 to 71, and for sweet spot 66 to 69. The smallest gap, three points for sweet spot, fits my experience of doing mathematics. I have a fairly good idea of what I'm good at and what work I'd willingly spend an afternoon on even if nobody had assigned it. The largest gap, twenty-two points for integrated life, is harder to sit with. If research, teaching, and club work each give me a perfectly good reason to spend an afternoon, I still have only the one afternoon. Naming a value that justifies each commitment does very little to resolve the conflict between them. The thirteen-point gap for support team also gives me something specific to consider before next year. If another member wants to organize an event, do I trust them enough to give them a part of the decision that I genuinely care about? That seems a fairer test of my willingness to share the work than counting how many people I know who would be willing to help.

\subsection{Active Experimentation and Closing Ethical Reflection}
In recitation I already tried allowing the people affected by a decision to change it. A slight majority of my students preferred projected notes, I switched, and a few thanked me. For the club I want to try a stronger version of that experiment. Before settling on a program, I'd like to ask members what they would come to even if they didn't know anybody there, and invite someone interested in an idea to decide with me how to run it. An idea I wouldn't have chosen making it onto the calendar would tell me something that a general assurance that everybody feels welcome couldn't. Servant leadership seems quite compatible with being the person who books the room. In fact, booking the room for somebody else's idea strikes me as a rather good reason to take the job.
\entry{Week 10. Humility \textit{(October 27 and 29, 2026)}}
\subsection{Concrete Experience}
I'm working with two professors in the math department, Dr. David Singer and Dr. Ryan Matzke on energy optimization and potential theory. One motivation for our project is the Thomson problem. If a fixed number of identical charges can move on the surface of a sphere, how should they arrange themselves to minimize their total electrostatic energy? One question I want to understand is what happens when we change what it means for two particles to be close together. Imagine two points on opposite sides of a sphere. I can measure the straight line joining them through the interior, or I can measure the shortest route between them along the surface. The first gives the chordal distance and the second the geodesic distance. In a Riesz energy, each pair contributes an inverse power of its distance, so the choice enters into the very quantity we are trying to minimize. We're interested in how results for the sphere extend to other surfaces, among them the Clifford torus, which one can think of as a product of two circles. There are several distinct questions even after we've fixed the surface and the distance. How close can two particles in an optimal configuration get to one another? As we allow more and more particles, how is their charge distributed over the surface as a whole? A description of the limiting distribution can tell us a great deal about the second question while still leaving open where to put a particular number of points.

I've spent weeks reading Borodachov, Hardin, and Saff's monograph on discrete energy with Dr. Matzke, the standard reference for the subject (Hardin was Dr. Matzke's PhD advisor!). Part of the work of reading it is learning to recognize which assumptions do the work in a proof. A theorem proved with straight-line distance comes with no guarantee that replacing that distance by shortest-path distance will preserve the argument. In my scholarship application I wrote that I expected to play a role in determining the direction of this project. I meant it when I wrote it. Apparently I could anticipate that responsibility in the abstract and still be surprised the first time it took the very ordinary form of two people wanting my help with a paper.

Recently the three of us were reading a paper together. I'm an undergraduate working with two professors, so I'm accustomed to there being a fairly straightforward answer to the question of who is helping whom. What took me by surprise was realizing that they were looking to me for guidance in understanding the paper. They wanted to know what I made of it. It occurred to me that they regarded me as an equal participant in our research, a possibility to which I apparently hadn't given very serious consideration until they had occasion to need my help.

\subsection{Reflective Observation}
There are at least two reasons a professor might ask a student how an argument works. One is that the professor knows the answer and would like to see whether the student can arrive at it. I do this myself in calculus recitation. The other is that the professor would like to understand the argument. These questions can have exactly the same wording, yet the difference matters a great deal to the person answering. In the first case I might reasonably spend some time figuring out what my professor is getting at. In the second, time spent guessing what answer Dr. Singer or Dr. Matzke wants is time that three people could have spent figuring out whether the argument works. I had carried a perfectly reasonable assumption from the classroom into a setting where it was beginning to get in the way.

\subsection{Abstract Conceptualization and Reading Reflection}
The Society of Leadership Fellows calls confidence one of the conditions for practicing humility (2020, pp. 35--36). At first that struck me as an odd qualification. Surely a person who is already inclined to think highly of their abilities has enough encouragement to do so. I have to consider, though, what the two professors actually risked. Either could acknowledge that a passage was giving them trouble without much danger of anybody concluding that they had no business doing mathematics. Their confidence gave them the freedom to ask an ordinary question. The more I worried about establishing my right to be there, the less free I was to give them an ordinary answer.

Kissinger's (2022, pp. 6--7) account of Adenauer supplies a hard case for any definition of humility that amounts to knowing one's place. As mayor of Cologne, Adenauer refused to welcome Hitler at the airport and had Nazi flags removed from public sites. He subsequently lost his office. I find the order of those events important. Had he understood humility as a standing obligation to yield to the person with greater authority, he could scarcely have acted as he did. The stakes in a research meeting are enormously smaller, but the difficulty with the definition is precisely the same. A title can tell me whose judgment merits considerable attention, but it can't settle on its own whether a particular judgment is correct.

\subsection{Assessment Reflection}
I scored 21 out of 30 for understanding myself, 25 for relating to others, and 27 for perspective. All three fall in the high range of 21--30. What interests me is the difference between two items that initially sounded almost interchangeable. I gave myself a five on being open to new information about myself and a three on actively seeking information about my weaknesses. I can quite sincerely accept a correction once somebody offers it while doing relatively little to make sure they have the opportunity. My sixes on the items about others' contributions, my place in a larger community, and my connection to a larger purpose also make sense of the meeting. I have no difficulty recognizing the two professors' contributions to our research. Apparently it takes more work to recognize my own place in that same community when the contribution happens to be mine. My six on enjoying cooperation gives me a practical reason to work on this. I enjoy doing mathematics with people, and we're more likely to get to do it together if I tell them where I think the argument works and invite them to show me where it fails.

\subsection{Active Experimentation and Closing Ethical Reflection}
For our next paper I'd like to bring one passage I can explain and one equation I can't quite justify. For the first I'll try giving the argument directly. For the second I'll identify the exact step at which I get lost, so that either professor can take it up there. The outcome I care about is whether one of us can carry the other's reasoning a step further. This also gives me a way of recognizing humility in my own students. If a student tells me my explanation skipped something, I want to be as willing to let them point it out as these two professors were to let me help. Of the things they could have taught me simply by knowing more mathematics, this particular lesson depended on their being willing to ask.

\entry{Week 11. Courage \textit{(November 3 and 5, 2026)}}
\subsection{Concrete Experience}
A number of my friends work on the student theater company, Footlighters, in one capacity or another.
After attending their production of the musical \textit{Alice by Heart}, I made the decision that I missed theater enough to try to get back into it.
Despite almost bailing out, I auditioned for the theater department's production of \textit{Company}. This was my first time doing any theater in four years. From sixth grade through twelfth I spent half of every school day in musical theater, so there is something a little ridiculous about how much courage it took to return to an activity on which I'd once spent that much time. If somebody had asked me in high school what I did, theater would have been a perfectly ordinary part of my answer. Four years later, I had to decide whether I was willing to give that answer again in front of people who could actually see me perform.

\subsection{Reflective Observation}
One might think that having done theater for years should make an audition easier. In one sense, of course, it does. I know what an audition is and what it means to commit to a performance. But there's a peculiar comfort to being a genuine beginner. Nobody, myself included, has a particularly well-developed idea of how well I ought to be able to do something I've never done before. Here I had a whole education's worth of reasons to expect something of myself, and four years' worth of reasons to wonder whether I could still do it. I could have continued to think of myself as somebody who used to be pretty good at theater indefinitely, and to be honest I really did. Auditioning meant being willing to find out something about the person who showed up this time.
A big reason why I forced myself to actually buck up and prepare for the audition was to stop myself from becoming one of those ``I could've gone pro in football if I hadn't hurt my knee,'' types.

\subsection{Abstract Conceptualization and Reading Reflection}
Ladkin's exercise on exemplars of moral courage (2021, p. 141) asks whom I admire, what that person did, and how I might emulate them. Earlier in this course I had a great deal of difficulty naming a single person whose character I could endorse without qualification. This time I'd choose an action. As Kissinger recounts, Adenauer used an authority he already possessed as mayor to refuse a public show of welcome to Hitler. What I admire is his recognition that even the apparently ceremonial business of greeting a visitor could communicate his city's assent. The question I can carry into a room with people who outrank me is what they might reasonably take my silence to mean. I have considerably less to lose than Adenauer did. I also know from my assessment that the possibility of displeasing a supervisor weighs on me.

Detert's (2022) account of Fred Keller offers a different example. Cascade Engineering's first two attempts to bring people out of poverty and into employment failed. One detail I particularly like is that Keller subsequently had his own managers, himself included, learn more about poverty. Repeating the hiring program indefinitely would have demonstrated considerable persistence while asking the next group of employees to bear the cost of the same mistakes. Keller was willing to expose his own ignorance to the people whose cooperation he needed. Going back to theater allowed a director to see what four years away had done to my performance. Detert leaves me with the harder question of whether I'd be equally willing to let a criticism change how I did something once other people depended on me.

One last thing I'll say is that there's a part of me that auditioned because I miss the arts, and a part of me that auditioned because I'm terrified about the future of academia and science, and I don't want to find myself deserted without the most human activity of all.
The advent of AI in academia scares me a lot, and I worry a ton about how replacable I am.
The one place I know for sure I'm not replaceable is in the arts.

\subsection{Assessment Reflection}
My scores on the professional moral courage assessment, averaging the three items in each theme, were 5.67 for moral agency, 4.67 for multiple values, 5.00 for endurance of threats, 6.00 for going beyond compliance, and 6.00 for moral goals. My fifteen answers totaled 82, for an overall mean of 5.47 out of 7. The two sixes make sense to me. I like to understand the purpose of a rule, and I'm fairly insistent that an action have a defensible end. The fours deserve more attention. In particular, I gave myself a four on acting morally when it puts me in an uncomfortable position with a supervisor. My affection for a good argument doesn't establish how readily I'll give one to a person who has power over me.

Would I call auditioning an act of moral courage? I think that would be generous. Nobody's welfare depended on my audition for \textit{Company}, and I didn't have to oppose an injustice to do it. What the two cases share is the experience of putting a judgment I care about in the hands of someone whose response I can't control. My lower score on multiple values suggests the point at which the analogy gives out. In a genuine moral disagreement, the other person may have a legitimate obligation that I haven't even considered. Showing up, difficult as it can be, settles very little about whose judgment is correct.

\subsection{Active Experimentation and Closing Ethical Reflection}
The class prompt about whether everybody else's conduct makes my own acceptable reminded me of the chalkboard in Strosacker. Teaching calculus at a board is ordinary enough that I could have carried on indefinitely with the reassurance that lots of good instructors do it. The fact that a practice works elsewhere says very little about whether the people in the last row of my room can see. My own interest in practicing board work made the appeal to convention especially convenient. I want to remember that the next time I catch myself using what most people do as a reason to stop asking questions.

Auditioning was itself the first experiment. Whatever the casting outcome, I now know that a four-year absence is compatible with going back. In research I can try a more specifically ethical application the next time I believe a step in a shared argument is mistaken. All three of us have an interest in getting the result right. I can identify the step, give my reason, and ask what assumption I'm missing. If I come away persuaded that I was wrong, the experiment may still have gone quite well. What I want to find out is whether concern about how I look prevents the group from examining a question that bears on its work.

Note added Wednesday, November 18th: I got a callback for Bobby (the leading man)! I'm super excited for the chance to sing some of my favorite music in the world, and I'll be happy with this regardless of the overall outcome.
\entry{Week 12. Accountability \textit{(November 10 and 12, 2026)}}
\subsection{Concrete Experience}
Earlier in the semester I asked my calculus students whether they preferred the chalkboard or the document camera. A slight majority preferred the document camera. I changed the way I taught recitation, and a few students stayed after class to thank me. In my first journal I wrote about this as an exercise in listening. What strikes me about it now is a detail that I hardly thought worth including the first time. The students could walk into Strosacker the following Thursday and see what I had done about their answers. They didn't have to infer from the warmth of my response whether I had taken them seriously. The projector was on.

\subsection{Reflective Observation}
I find it relatively easy to take responsibility for a decision when I also get to tell the story of how well I handled it. Imagine that I'd collected the same responses, written a thoughtful journal entry about my surprise, and returned to the board the following week. From my perspective a great deal would have occurred. I'd have questioned my motives, considered the competing reasons, and perhaps even learned something about myself. From a student's perspective I would have asked a question and carried on as before. They could reasonably have wondered what I wanted their feedback for. Much of what I count as the work of being conscientious is invisible to precisely the people who have an interest in knowing what I'll do.

\subsection{Abstract Conceptualization and Reading Reflection}
In \textit{Leading with Integrity}, the Society of Leadership Fellows (2019b, ch. 3) discusses the questions that make it harder for a questionable practice to persist unnoticed. I had thought of my feedback form primarily as a way to let students say what they preferred. The more uncomfortable question was one a student could quite reasonably have put to me. Whose education was being served by my practicing at the chalkboard? I had a good answer for why learning to lecture at a board mattered to me. It was the ease with which I could supply that answer that had kept me from asking the second question about the people sitting toward the back of Strosacker.

Rea, Stoller, and Kolp (2023, ch. 13) suggest looking for the forces that keep an existing habit in place. In my case the room had an enormous board, I wanted practice, and a student who was struggling to see had to make the effort to tell the very person who had chosen the format. The survey changed the last condition. Thinking about the November 12 community-of-wise-others exercise, I can see a limit to my usual choice of advisers. My father has taught me a great deal about how to treat people. He can't tell me how legible my handwriting looks from the last row of a particular auditorium. Sometimes the person whose judgment I need is simply the person sitting there.

\subsection{Assessment Reflection}
I scored 67 out of 100 for clear ownership, 67 for visible follow-through, and 69 for constructive response. Two points is a fairly slender basis on which to claim that constructive response is my great strength. What I do recognize in the distinction is that I like the moment somebody brings me a problem. There is something interesting to think about, and very often there is an immediate opportunity to be helpful. The person may have a wholly different interest in the matter. To a student waiting to find out how recitation will be taught, my initial willingness to engage might be the least consequential part of the interaction. The near equality of the scores gives me a reason to attend to all three stages of a commitment, particularly the less intellectually interesting business of making sure the person who asked knows what happened.

\subsection{Active Experimentation and Closing Ethical Reflection}
The accountability worksheet asks what agreements a new project team needs at the outset and how a leader ought to respond to a mistake. For the next club event I'd like each of us to be able to say who has agreed to contact a speaker, by when, and who needs to hear back. If I've misunderstood a conversation and nobody has sent the invitation, my first job is to tell the group what happened and establish who can still reach the speaker in time. Working out how the misunderstanding arose will be easier once the event is no longer getting later by the minute.

I recently had a chance to practice this kind of follow-through on a larger scale. I organized a professional development day for the Association for Women in Mathematics (AWM). We brought in a graduate student in applied mathematics, a professor in pure mathematics, and an alum working in actuarial science to talk with undergraduates. Getting all three there took real work. A date that suited one panelist was of little use until I knew the other two could make it, and the department needed enough advance notice to cater the event. Planning that far ahead also gave undergraduates time to set the day aside. We had a large turnout, the math department catered, and the department chair came. I felt an enormous sense of accomplishment seeing those three people together in front of the students. Earlier in the semester I might have described my part as merely scheduling things. Having been responsible for getting everybody there, I find it difficult to say what would have been left of the panel without somebody doing precisely that work.

There was a reason to go to the trouble of coordinating those three schedules. As somebody who wants to pursue research, I can say considerably more about what interests me in pure mathematics than I can about what an actuary does with a mathematics degree. Had I settled for the easiest panel for me personally to organize, my own familiarity with one path could easily have passed for a fair account of the choices available to an undergraduate. With the three panelists there, a student could consider graduate study in applied mathematics, academic work in pure mathematics, and a career outside the university in the same sitting. It also mattered when I could tell people about it. An invitation to undergraduates is worth rather less if it comes after they've already committed the afternoon elsewhere. The same was true of asking the department to provide food. The time I spent planning in advance gave other people time to decide what they could contribute. That's one sense of respect I hadn't thought much about before this class.

I also want to be precise about what the day validated for me. The department chair's presence meant something to me, and I was proud to see a large audience. But if I needed the chair to attend before I could regard the work as worthwhile, I'd have much the same problem I described in the authenticity entry. What I can take credit for is having followed through on a series of decisions whose consequences were visible to other people. Next year, when I worry about taking up too much space in the math club, I can ask a question for which I now have an actual example. Whose opportunity to speak or participate might depend on my willingness to do the organizing?

\entry{Week 13. Self-Compassion \textit{(November 17 and 19, 2026)}}
\subsection{Concrete Experience}
About a week and a half ago I realized how much I was struggling with anxiety and self-control, and how long it had been going on. Looking back, I could see the pattern over a number of months. I reached out to my doctor and made an appointment. I told her that some of my behavior reminded me of myself in high school, in ways I had hoped I'd grown out of. She commended me on the self-awareness and maturity it took to recognize that in myself. We talked through a number of steps I could take, and I've begun working on the habits I'd fallen into over the past few months.

\subsection{Reflective Observation}
What interested me about my doctor's response was that we were considering the same piece of evidence and drawing conclusions about two different things. I recognized something I'd done in high school and took the resemblance as a reason to worry about who I was becoming. She could also see a person who had noticed the resemblance, called her, and come prepared to talk about it. Both accounts are true. For months I had given most of my attention to my work and to other people. I assumed that if there was a cost to putting myself last, I would be the one to bear it. By the time I made the appointment, the anxiety and difficulty with self-control were affecting the work to which I'd been so determined to give my attention.

\subsection{Abstract Conceptualization and Reading Reflection}
In the self-care exercise, Hillel asks in succession what happens if I leave responsibility for myself to others, what I become if I care only for myself, and when I intend to act (Pirkei Avot 1:14). I'd been taking the second question seriously enough that I'd given myself a fairly convenient way to evade the first and third. There was always somebody else I could help that day. Eisler (2017) makes a related point about what an accounting system renders visible. A national account can record what it costs to clean up an oil spill while leaving out much of the care that households provide every day. My own accounting had a similar oddity. I could tell you what work I'd gotten done at considerable cost to myself, but it was only when that cost began to affect my work that I took it seriously enough to do something about it. I'd like to think my own welfare gives me sufficient reason somewhat earlier than that.

Rea, Stoller, and Kolp (2023, ch. 14) distinguish what a number can establish from the meaning a story can convey. My 1.80 for self-kindness is recognizable enough, but the score alone could never have told me what I heard in my own entries. Over time, the voice in which I was describing myself had grown more critical. Journaling gave me some distance from thoughts that otherwise sounded quite reasonable simply because I was the one thinking them. I think having that record helped me recognize that I needed a little help.

\subsection{Assessment Reflection}
My overall score was 2.65 out of 5, within the report's moderate range and below its approximate average of 3. On the positive subscales, I scored 1.80 for self-kindness, 3.00 for common humanity, and 3.25 for mindfulness. On the negative subscales, I scored 3.40 for self-judgment, 3.75 for isolation, and 3.00 for over-identification. Higher raw scores on the negative subscales indicate more frequent negative responses. Isolation was the highest of those three, which seems odd given how readily my mother and my professors helped with my Goldwater application. Then again, the question of whether there are people who would help me is quite distinct from the question of whether, at the time I'm struggling, I regard my struggle as something those people might understand. My parents could recognize how much I cared about the work and how much it was costing me in the same conversation. I had given myself considerably less latitude.

\subsection{Active Experimentation and Closing Ethical Reflection}
Making the appointment was a first experiment in responding differently to a familiar pattern. My doctor gave me specific steps, and I have begun taking them. It is early enough that the most honest outcome I can report is that I've started. In future entries I'd like to pay attention to the first sentence I write about a difficult day. Did I describe something I did that I'd like to change, or did I go straight from one bad afternoon to an assessment of my entire character? If I can catch that jump on paper, I have a better chance of doing something about the difficulty that prompted it.

\entry{Synthesis of Ethical Leadership Learning}
When I wrote in October that I'd rather be someone the Guy could depend on than be the Guy myself, I had a fairly specific person in mind, even if I hadn't given him a name. He was the one people looked to when they needed guidance. My job, as I imagined it, was to know enough and care enough to help him give it. It hadn't occurred to me to ask what would happen if the person I considered the Guy turned around and asked me how a proof worked. Meeting with Dr. Singer and Dr. Matzke made the question considerably less hypothetical. They were willing to do exactly what I admired in a colleague. They wanted help understanding something, and they asked the person they thought could help. It was my own account of my place in the room that made their conduct difficult to understand.

I think my worry about leading the math club has something in common with that surprise. I can say fairly confidently what I want members to have. I want them to be able to come with a question, to spend time with people who like thinking about it, and to leave feeling that they had some say in what the club does. I have much less confidence in my judgment of how conspicuous I can be before I'm taking up too much space. Fortunately, these questions admit of different kinds of evidence. Trying to gauge how conspicuous I feel is an exercise at which I could spend an almost unlimited amount of time. Asking whose questions end up on the program, or whether a new member has ever had an idea survive my initial lack of enthusiasm, might yield an answer of some consequence to another person. Our earlier class discussion about whether an organization can care and Gorman's way of bringing a particular person's voice into a collective project give me a way to understand servant leadership that I'd be willing to practice. The members need some real authority over what counts as serving them.

The audition complicates this account in a way I like. I didn't audition for \textit{Company} in order to serve anybody, I just wanted to do theater again. If I'm too quick to find a lofty ethical purpose in the decision, I miss what was difficult about it. I had to want something enough to give other people the opportunity to tell me I might not get it. The moral courage assessment gives me a reason to be careful about how far I generalize from that experience. I scored highly on thinking beyond compliance and considering moral ends, and somewhat lower on enduring threats and reconciling different values. It's entirely possible to walk into an audition and still keep quiet about an error in front of a professor. What I can take from the audition is more modest and more credible. Waiting to feel as sure of myself as I did four years ago needn't be a condition of doing something I have reason to do now.

Dr. Case wanted me to be specific about how the course has helped me become a better leader. The AWM professional development day is the clearest answer I have. In my first journal I described myself as someone who wanted to be useful to the person in charge. For the panel, I could name something the undergraduates deserved to have and take responsibility for making it happen. Thinking in terms of accountability changed what counted as having done my part. Sending invitations wouldn't give students the opportunity to compare applied research, pure mathematics, and actuarial work. For that, I needed three people in the same room at a time the students knew about far enough in advance to come. Seeing the turnout, the panelists, and the department chair there made me proud of having done the coordinating. It also gives me a reason to trust myself with the math club next year that I didn't have when all I could offer was a description of the leader I hoped to be.

When Dr. Case and I met on Tuesday, November 17, we discussed how much trouble I'd had with reflexive journaling during the first seven weeks. I think I've finally figured out what I get out of doing it. On any given day I can give a fairly convincing explanation of why I'm anxious, why one more piece of work needs my attention, or why I'm disappointed in how I handled something. What I have trouble hearing from inside any one of those explanations is how critical the voice giving it has become. Once writing became a habit, I could follow that voice across entries. I think that helped me recognize the need to talk to my doctor. It also changed what I understand self-awareness to require. I used to think the work was largely done once I could give a sufficiently accurate account of myself. This semester one of the more mature things I did with such an account was to take it into somebody else's office.

One last thing I'll say is this has been a uniquely strange semester for me.
I came into it knowing it was going to be tough, what with all of my classes, club commitments, research projects, etc., and I thought I was prepared.
But when I started to notice myself falling back into old habits, I got scared.
I remembered how I felt in high school, when I was failing my way through classes, just surviving until it was time to go to rehearsal or music class, and I psyched myself out.
With the support of friends, my parents, and my doctor, though, I think I'm moving in the right direction now, and I genuinely think this journal and this class are a small part of that progress.
I've tried to keep a journal before, and never really succeeded to be consistent.
Ironically, having a grade attached to it was the one thing that actually got me to sit down and write my thoughts every week.

\entry{References}
Detert, J. R. (2022, January 7). What courageous leaders do differently. \textit{Harvard Business Review}.

Eisler, R. (2017). The real wealth of nations: From global warming to global partnership. \textit{Interdisciplinary Journal of Partnership Studies, 4}(3), Article 13.

Gorman, A. (2021, January 20). \textit{The hill we climb} [Inaugural poem].

Borodachov, S. V., Hardin, D. P., \& Saff, E. B. (2019). \textit{Discrete energy on rectifiable sets}. Springer.

King, M. L., Jr. (1963, August 28). \textit{I have a dream} [Speech]. March on Washington for Jobs and Freedom.

Kissinger, H. (2022). Konrad Adenauer: The strategy of humility. In \textit{Leadership: Six studies in world strategy} (pp. 3--48). Penguin Press.

Ladkin, D. (2021). Taking ethical action. In \textit{Mastering ethics in organizations: A self-reflective guide to developing ethical astuteness} (ch. 7). Edward Elgar.

Newman, A. (2019). \textit{Building leadership character}. SAGE.

Rea, P., Stoller, J. K., \& Kolp, A. (2023). Forming good habits and The value of virtue. In \textit{Better humans, better performance} (chs. 13--14). McGraw Hill.

Self-Compassion Institute. (2026, November 18). \textit{Your personal self-compassion scores} [Individual assessment report].

Society of Leadership Fellows. (2019a). \textit{Authenticity and trust in leadership: Insights from three leadership conversations}. College of St. George.

Society of Leadership Fellows. (2019b). \textit{Leading with integrity: Insights from three leadership conversations}. College of St. George.

Society of Leadership Fellows. (2020). \textit{Courage and humility in leadership: Insights from five leadership conversations}. College of St. George.
\end{document}
